\appendix
\section{Statistical Analysis}
\label{app:stats}

\Cref{tab:stats} reports per-seed ASR and Utility (mean over 3 seeds, 30 task-injection pairs per seed) with 95\% Wilson confidence intervals, computed from the native end-to-end results (990 main trials + 270 ablation trials). All values are from AgentDojo's \texttt{security\_from\_traces} (ASR) and \texttt{user\_task.utility} (Utility) evaluators on real tool-execution traces.

\begin{table}[h]
\centering
\small
\caption{Per-defense end-to-end results: ASR and Utility with 95\% Wilson CI (n=90 per defense, 3 seeds).}
\label{tab:stats}
\begin{tabular}{lcc}
\toprule
Defense & ASR [95\% CI] & Utility [95\% CI] \\
\midrule
No Defense & 0.74 [0.65, 0.82] & 0.19 [0.12, 0.28] \\
Polymorphic & 0.14 [0.09, 0.23] & 0.73 [0.63, 0.81] \\
Mixture-of-Enc. & 0.16 [0.10, 0.24] & 0.69 [0.59, 0.78] \\
Spotlighting & 0.21 [0.14, 0.31] & 0.67 [0.57, 0.76] \\
StruQ (prompt) & 0.70 [0.60, 0.78] & 0.17 [0.10, 0.27] \\
FATH & 0.00 [0.00, 0.04] & 0.02 [0.00, 0.07] \\
Task Shield & 0.10 [0.05, 0.18] & 0.02 [0.00, 0.07] \\
IPIGuard & 0.38 [0.28, 0.48] & 0.06 [0.02, 0.13] \\
P1: Echo Canary & 0.03 [0.01, 0.09] & 0.00 [0.00, 0.04] \\
P2: Orthogonal & 0.00 [0.00, 0.04] & 0.00 [0.00, 0.04] \\
P3: Task Invariant & 0.42 [0.32, 0.52] & 0.02 [0.00, 0.07] \\
\bottomrule
\end{tabular}
\end{table}

\section{P2 Ablation Detail}
\label{app:p2_ablation}

\Cref{tab:p2_ablation_app} details the P2 ablation (90 trials per configuration, same tasks/attacks/seeds as the main experiment). The canary provides independent security value (ASR 0.03 alone vs.\ 0.74 baseline), but all configurations sacrifice Utility entirely.

\begin{table}[h]
\centering
\small
\caption{P2 ablation: canary vs.\ alignment contributions (n=90 each).}
\label{tab:p2_ablation_app}
\begin{tabular}{lcc}
\toprule
Configuration & ASR [95\% CI] & Utility [95\% CI] \\
\midrule
Canary only & 0.03 [0.01, 0.09] & 0.00 [0.00, 0.04] \\
Alignment only & 0.09 [0.05, 0.17] & 0.04 [0.02, 0.11] \\
Combined (P2) & 0.00 [0.00, 0.04] & 0.00 [0.00, 0.04] \\
\bottomrule
\end{tabular}
\end{table}

\section{Over-Blocking Diagnosis}
\label{app:overblock}

To distinguish genuine defense over-blocking from base-model capability ceiling, we compare Task Shield to the no-defense baseline on identical (seed, user\_task, injection\_task) triples. Of 90 triples: the baseline completes 13 (Utility=1); Task Shield blocks all 13 of these (genuine over-blocking = 14.4\%). Both fail on 72 triples (model weakness = 80.0\%). Task Shield completes 5 triples the baseline fails (5.6\%). This analysis shows that the low Utility of runtime-checking defenses is dominated by the base model's capability ceiling, not by defense over-blocking---but genuine over-blocking does exist (14.4\% of completable tasks blocked).

\section{Omitted Prior Experiments}
\label{app:omitted}

The prior version of this paper (proxy-based evaluation) included backbone comparisons (GPT-4o, Qwen2.5-7B), cross-benchmark validation (InjecAgent), adaptive-attack-strength sweeps, and SecAlign evaluation. These were conducted under the string-matching proxy methodology and produced results that \emph{contradict} our native end-to-end findings (e.g., proxy reported Spotlighting R-ASR=1.00; native shows ASR=0.21). We have therefore \emph{removed} these proxy-based results from this revision rather than report numbers we cannot support with end-to-end evidence. We plan to re-run backbone, cross-benchmark, and SecAlign evaluations using the native pipeline in future work.

\section{Native Evaluation Pipeline}
\label{app:native}

Our evaluation uses AgentDojo v1.2.2's native benchmark infrastructure. Each defense is implemented as a \texttt{BasePipelineElement} that integrates into AgentDojo's \texttt{AgentPipeline} and \texttt{ToolsExecutionLoop}. The agent executes real tools (e.g., \texttt{send\_email}, \texttt{get\_calendar\_events}), producing real function-call traces and environment state changes. Security is scored by \texttt{injection\_task.security\_from\_traces(traces, pre\_env, post\_env)} and utility by \texttt{user\_task.utility(model\_output, pre\_env, post\_env)}. This eliminates the string-matching proxy used in prior work. All code is in \texttt{sok\_ipl/native/pipeline\_builder.py} and \texttt{scripts/run\_native\_benchmark.py}.
